\documentclass[12pt,a4paper]{article}

\usepackage[a4paper, top=0.9in, bottom=0.9in, left=1in, right=1in]{geometry}
\usepackage{graphicx}
\usepackage{hyperref}
\usepackage{listings}
\usepackage{xcolor}
\usepackage{booktabs}
\usepackage{array}
\usepackage{titlesec}
\usepackage{fancyhdr}
\usepackage[T1]{fontenc}
\usepackage[utf8]{inputenc}
\usepackage{parskip}
\usepackage{enumitem}
\usepackage{microtype}
\usepackage{tabularx}
\usepackage{mdframed}
\usepackage{amssymb}
\usepackage{amsmath}

\hypersetup{colorlinks=true, linkcolor=blue!70!black, urlcolor=blue!70!black}

\definecolor{codebg}{rgb}{0.96,0.96,0.96}
\definecolor{codecomment}{rgb}{0.25,0.50,0.35}
\definecolor{codekeyword}{rgb}{0.00,0.00,0.65}
\definecolor{codestring}{rgb}{0.63,0.13,0.13}
\lstset{
    backgroundcolor=\color{codebg},
    basicstyle=\ttfamily\footnotesize,
    keywordstyle=\color{codekeyword}\bfseries,
    commentstyle=\color{codecomment}\itshape,
    stringstyle=\color{codestring},
    numbers=left,
    numberstyle=\scriptsize\color{gray},
    stepnumber=1,
    numbersep=6pt,
    frame=single,
    rulecolor=\color{gray!50},
    breaklines=true,
    captionpos=b,
    tabsize=4,
    showstringspaces=false,
    xleftmargin=1.6em,
    framexleftmargin=1.2em,
    aboveskip=8pt,
    belowskip=6pt,
}

\titleformat{\section}{\large\bfseries\color{blue!75!black}}{}{0em}
    {\thesection\quad}[\vspace{-2pt}\rule{\linewidth}{0.5pt}]
\titleformat{\subsection}{\normalsize\bfseries}{}{0em}{\thesubsection\quad}
\titlespacing*{\section}{0pt}{14pt}{6pt}
\titlespacing*{\subsection}{0pt}{10pt}{4pt}

\newmdenv[
    backgroundcolor=blue!4,
    linecolor=blue!35,
    linewidth=0.8pt,
    innerleftmargin=8pt,
    innerrightmargin=8pt,
    innertopmargin=6pt,
    innerbottommargin=6pt,
    skipabove=8pt,
    skipbelow=8pt,
]{grammarbox}

\pagestyle{fancy}
\fancyhf{}
\rhead{\small MuleTrace (MTX) --- Real-Time AML Decision Engine}
\lhead{\small IIITDM Kurnool}
\cfoot{\thepage}
\renewcommand{\headrulewidth}{0.4pt}

\setlength{\parskip}{6pt}
\setlength{\parindent}{0pt}

\begin{document}

% ── PAGE 1: Title ─────────────────────────────────────────────────────────────
\begin{titlepage}
    \centering
    \vspace*{0.8cm}
    {\large\bfseries\color{blue!75!black}
        Indian Institute of Information Technology,\\[2pt]
        Design and Manufacturing, Kurnool}\par
    \vspace{1.4cm}
    {\LARGE\bfseries AML Money-Mule Platform Report}\par
    \vspace{0.4cm}
    {\normalsize December 2025 to June 2026}\par
    \vspace{1.8cm}
    {\LARGE\bfseries MuleTrace (MTX)}\par
    \vspace{0.3cm}
    {\large\itshape A Real-Time AML Money-Mule Classification \& Entity Resolution Platform}\par
    \vspace{1.8cm}
    {\normalsize Indian Institute of Information Technology,\\
    Design and Manufacturing, Kurnool}\par
    \vspace{1.6cm}
    {\normalsize Submitted by:}\par\vspace{0.4cm}
    {\large\bfseries M.V.S Karthikeya}\par
    {\normalsize Roll No:\ 124CS0021}\par
    \vspace{0.5cm}
    {\large\bfseries R. Shailesh}\par
    {\normalsize Roll No:\ 524CS0010}
\end{titlepage}

% ── PAGE 2: Project Overview ──────────────────────────────────────────────────
\section*{Project Overview}
\addcontentsline{toc}{section}{Project Overview}

\textbf{MuleTrace (MTX)} is a service-oriented, low-latency Anti-Money Laundering (AML) classification and entity resolution platform. It combines gradient-boosted machine learning inference with policy-driven risk scorecards to classify bank accounts and discover hidden syndicates. The platform targets money mule networks, which route illicit funds using vulnerable customer demographics.

MuleTrace deploys a bimodal machine learning pipeline (XGBoost/LightGBM) to balance alert precision and recall, paired with a Local Feature Attribution (LFA) Explainable AI (XAI) engine. To comply with banking regulation (such as the US Federal Reserve Board's SR 11-7 Model Risk Management guidelines), the system automatically translates feature contributions into natural language compliance narratives. It also hosts an interactive auditing sandbox for manual scenario perturbation.

\subsection*{Project Objectives}
\begin{itemize}[itemsep=3pt]
    \item Implement a bimodal classification pipeline using LightGBM and XGBoost.
    \item Address extreme class imbalance where money mules represent only 0.89\% of cases.
    \item Mitigate systemic training data leakages (sorting and temporal batch leakages).
    \item Build a FastAPI REST service to support low-latency transaction scoring.
    \item Design a perturbation-based scorecard engine to audit manual policy overrides.
    \item Graph network linkages dynamically based on shared metadata (branch and occupation).
    \item Translate local machine learning feature importances into readable forensic explanations.
\end{itemize}

\subsection*{Technology Stack}
\begin{itemize}[itemsep=3pt]
    \item \textbf{Backend Web Services:} FastAPI, Uvicorn, Python 3.10
    \item \textbf{Machine Learning:} XGBoost, LightGBM, Scikit-learn, NumPy, Pandas
    \item \textbf{Frontend UI:} HTML5, CSS3, Vanilla JavaScript, TailwindCSS, FontAwesome
    \item \textbf{Data Visualization:} vis.js (dynamic interactive network graphs)
    \item \textbf{Infrastructure & Deployment:} Render Web Services (PaaS)
\end{itemize}

\subsection*{Platform Components}
\begin{enumerate}[itemsep=2pt]
    \item Ingestion & Preprocessing \hfill \texttt{app.py}, \texttt{model\_metadata.json}
    \item Bimodal ML Inference Pipeline \hfill \texttt{xgb\_hybrid\_model.json}, \texttt{lgb\_hybrid\_model.txt}
    \item Perturbation-Based Scorecard Engine \hfill \texttt{app.py} (predict\_account)
    \item Local Feature Attribution (XAI) \hfill \texttt{app.py}, \texttt{medians}
    \item Natural Language Compliance Generator \hfill \texttt{static/index.html} (generateNlpExplanation)
    \item Entity Resolution Network Graph \hfill \texttt{static/index.html} (vis.js integration)
\end{enumerate}

% ── PAGE 3: Table of Contents ─────────────────────────────────────────────────
\newpage
\tableofcontents

% ── PAGE 4: Dataset Architecture ──────────────────────────────────────────────
\newpage
\section{Dataset Architecture \& Feature Mappings}

\subsection{Raw Data Characteristics}
The raw training dataset contains high-dimensional transaction activity summaries. It exhibits an extreme class imbalance, where confirmed money mule accounts constitute only \textbf{0.89\%} of the dataset. The total database contains \textbf{3,921 raw features} representing account behavioral metrics, location histories, and device logins.

\subsection{Anonymization Strategy}
To protect customer confidentiality and comply with data privacy laws, all 3,921 feature names were anonymized using sequential index codes (\texttt{F1} through \texttt{F3921}). 

To construct an explainable compliance platform, statistical correlation audits isolated the \textbf{21 most critical features} recommended by bank fraud experts. These features are mapped to their functional banking indicators in Table~\ref{table:mappings}. All remaining baseline features maintain their anonymized IDs to prevent model validation divergence.

\subsection{Feature Map Reference}
\renewcommand{\arraystretch}{1.15}
\begin{tabularx}{\textwidth}{>{\ttfamily\small}p{2.5cm} >{\small}p{4.2cm} X}
\toprule
\textbf{Anonymized ID} & \textbf{Semantic Label} & \textbf{Compliance & Auditing Context} \\
\midrule
F115 & Daily Limit Ratio & Percent of daily transaction limits utilized. \\
F321 & Tx Freq / Day & Outgoing transaction frequency per day. \\
F527 & Tx Activity A & Inbound money transfer volume. \\
F531 & Tx Activity B & Outbound money transfer volume. \\
F670 & Fraud Link Flag & direct network connection to a known fraud entity. \\
F1692 & Hi-Risk Transfer Count & Outgoing transfer count to blacklisted branches. \\
F2082 & Net Assoc Score A & Login association scoring for remote networks. \\
F2122 & Net Assoc Score B & KYC verification risk index. \\
F2582 & Velocity Score & Debit volume standard deviation over a 7-day window. \\
F2678 & Device Vel Index & Count of distinct device switches. \\
F2737 & Loc Change Speed & Geodetic login speed between transactions (km/h). \\
F2956 & Off-Hrs Logins & Frequency of logins between 00:00 and 04:00. \\
F3043 & Alert Freq Index & Number of historical flags generated. \\
F3811 & Avg Balance & 30-day moving average of account balance. \\
F3836 & Net Balance Ind. & Balance deviation relative to cohort medians. \\
F3886 & Account Type & Core product (Savings, Current, Gold, etc.). \\
F3887 & Bank Branch Code & Unique bank branch registry identifier. \\
F3889 & Active Window & Total duration of active history. \\
F3891 & Registered Occupation & Declared job status (student, housewife, retired, etc.). \\
F3894 & Customer Age & Customer age; detects youth recruitment rings. \\
F3898 & Mule Proxy Indicator & Proxy score determining sandbox baseline status. \\
\bottomrule
\end{tabularx}
\label{table:mappings}

% ── PAGE 5: Preprocessing & Data Leakage Mitigations ──────────────────────────
\newpage
\section{Preprocessing \& Data Leakage Resolution}

\subsection{Sorting and Index Leakage}
During initial model audits, we discovered a major data leakage: the training dataset was ordered sequentially by the target class (confirmed money mules clustered at the end). 
\begin{grammarbox}
\textbf{Index/Sorting Leakage Rule:}\\
A decision tree or gradient boosted model trained on un-shuffled sorted data will split on the implicit row index or internal database pointers. This results in an artificial validation accuracy of \textbf{99.9\%} that collapses to \textbf{50.0\%} (random guess) in a real-time production stream.
\end{grammarbox}

\textbf{Mitigation:} The pipeline explicitly drops row index columns, shuffles the entire dataset using a fixed random seed prior to train-test splits, and resets the index pointers. This forces the trees to split on transactional behaviors rather than dataset position.

\subsection{Temporal Batch Leakage}
Two specific columns, \texttt{F2230} and \texttt{F3912}, were found to contain exact timestamp records or batch processing dates. Because clean cases and money mule cases were collected in distinct calendar months, these features became perfect target separators.
\begin{grammarbox}
\textbf{Temporal Batch Leakage Rule:}\\
Including temporal batch variables allows the trees to partition the data by time intervals instead of behavioral characteristics. If a fraud ring is active in March, the model flags March transactions as fraudulent, failing to generalize to future transaction windows.
\end{grammarbox}

\textbf{Mitigation:} Features \texttt{F2230} and \texttt{F3912} were completely excluded from the training feature set. This guarantees the model remains time-invariant and focuses on behavioral patterns like limit ratios and transaction velocities.

% ── PAGE 6: Model Pipeline ────────────────────────────────────────────────────
\newpage
\section{Machine Learning Model Pipeline}

\subsection{Bimodal Dual-Engine Strategy}
Rather than deploying a single model, MuleTrace implements a bimodal dual-engine configuration. This architecture gives compliance officers the ability to select the optimal model based on operational volume and risk:
\begin{enumerate}
    \item \textbf{XGBoost (Precision Focus):} Tuned to minimize false positives. It is used during periods of high investigator backlog to ensure legitimate customer accounts are not locked.
    \item \textbf{LightGBM (Recall Focus):} Tuned to maximize detection rate. It is used during fraud surges to capture as many suspicious accounts as possible.
\end{enumerate}

\subsection{Class Imbalance Mitigation}
To counter the 0.89\% minority class distribution, we applied class weight scaling. The positive class weight parameter \texttt{scale\_pos\_weight} was set to \textbf{110.0}, penalizing misses on positive mule cases by 110x relative to clean cases. Overfitting was controlled by restricting tree depth (\texttt{max\_depth = 4}) and limiting the number of terminal leaves (\texttt{num\_leaves = 15}).

\subsection{Validation Performance}
Model evaluation metrics on hold-out validation partitions are summarized in Table~\ref{table:eval}. Both models use a default risk threshold of \textbf{0.50}.

\vspace{10pt}
\begin{table}[ht]
\caption{Model Performance Evaluation Summary}
\label{table:eval}
\centering
\begin{tabular}{lccccc}
\toprule
\textbf{Model Engine} & \textbf{Accuracy} & \textbf{Precision} & \textbf{Recall} & \textbf{F1-Score} & \textbf{PR-AUC} \\
\midrule
XGBoost (Precision) & 99.8\% & 100.0\% & 81.25\% & 0.897 & 0.940 \\
LightGBM (Recall) & 99.8\% & 92.86\% & 81.25\% & 0.867 & 0.881 \\
\bottomrule
\end{tabular}
\end{table}

The XGBoost model achieved 100.0\% Precision on the validation set, confirming that zero false-positive alerts were generated, while the LightGBM model prioritized sensitivity with a high PR-AUC of 0.881.

% ── PAGE 7: REST API ──────────────────────────────────────────────────────────
\newpage
\section{REST API Specifications}

The scoring engine is exposed via a FastAPI backend. The primary transactional scoring endpoint is \texttt{/predict}.

\subsection{Request Payload Schema}
The request requires the target profile features, baseline features (for sandbox calculations), the classification engine type, and the decision threshold.

\begin{lstlisting}[caption={JSON Request Payload --- /predict}]
{
  "features": {
    "account_id": "SBX-207",
    "F3886": "Savings",
    "F3891": "student",
    "F3887": 94,
    "F321": 1.000,
    "F3811": 50000,
    "F2582": 0.000,
    "F3894": 34,
    "F670": 1.0,
    "F3836": 381286
  },
  "base_features": {
    "account_id": "ACC06824",
    "F3886": "Savings",
    "F3891": "student",
    "F3887": 94,
    "F321": 1.000,
    "F3811": 50000,
    "F2582": 0.000,
    "F3894": 34,
    "F670": 0.0
  },
  "engine": "xgboost",
  "threshold": 0.50
}
\end{lstlisting}

\subsection{Response Payload Schema}
The response returns the probability score, classification decision, and the top risk drivers sorted by attribution.

\begin{lstlisting}[caption={JSON Response Payload --- /predict}]
{
  "account_id": "SBX-207",
  "engine": "xgboost",
  "probability": 0.6501,
  "prediction": 1,
  "flagged": true,
  "threshold": 0.50,
  "top_drivers": [
    {
      "feature": "F670",
      "value": 1.0,
      "importance": 0.0045,
      "contribution": 0.350,
      "impact": "High Risk Driver"
    },
    {
      "feature": "F3836",
      "value": 381286.0,
      "importance": 0.0022,
      "contribution": 0.300,
      "impact": "High Risk Driver"
    },
    {
      "feature": "F3811",
      "value": 50000.0,
      "importance": 0.0018,
      "contribution": -0.013,
      "impact": "Mitigating Factor"
    }
  ]
}
\end{lstlisting}

% ── PAGE 8: Scorecard Engine ──────────────────────────────────────────────────
\newpage
\section{Interactive Sandbox \& Scorecard Engine}

\subsection{Profile Cloning Mechanism}
Gradient-boosted trees require complete, non-null input vectors to perform inference. To evaluate hypothetical scenarios without manual entry of all 3,921 columns, MuleTrace uses a \textbf{Profile Cloning} mechanism. The investigator clones a real production account profile (which populates all baseline values) and modifies only a few key features. The backend merges these changes with the baseline vector and scores the result.

\subsection{Scorecard Adjustment Rules}
When an account ID begins with the prefix \texttt{SBX-}, the engine enters sandbox mode. It compares the perturbed values against their baselines and applies scorecard rules:

\begin{grammarbox}
\textbf{Risk Escalation Surcharges (Applied to Clean Profiles):}\\
- \textbf{Fraud Link flag (\texttt{F670} = 1.0):} Adds $+35.0\%$ to the risk score.\\
- \textbf{High-risk transfer count (\texttt{F1692} > 5.0):} Adds $+25.0\%$ to the risk score.\\
- \textbf{Demographic vulnerability (\texttt{F3894} < 23):} Adds $+10.0\%$ to the risk score.\\
- \textbf{Youth balance risk:} Student/housewife with balance (\texttt{F3811}) > \$1M adds $+50.0\%$; > \$250k adds $+30.0\%$; > \$50k adds $+15.0\%$.
\end{grammarbox}

\begin{grammarbox}
\textbf{Risk Mitigation Offsets (Applied to Suspect Profiles):}\\
- \textbf{Salaried occupation (\texttt{F3891} = "salaried"):} Deducts $-25.0\%$ risk.\\
- \textbf{Balance reduction (\texttt{F3811} < \$100k):} Deducts $-30.0\%$ risk.\\
- \textbf{Low transfer frequency (\texttt{F321} < 3.0):} Deducts $-20.0\%$ risk.\\
- \textbf{Disconnection from fraud network (\texttt{F670} = 0.0):} Deducts $-25.0\%$ risk.
\end{grammarbox}

\subsection{Reconciliation Integrity}
If an investigator clones a profile but submits it without making any modifications, the engine calculates a scorecard deviation of exactly \textbf{0.00\%}. This ensures the sandbox output reconciles perfectly with the production baseline.

% ── PAGE 9: Entity Resolution ─────────────────────────────────────────────────
\newpage
\section{Entity Resolution \& Graph Linkage}

\subsection{Metadata Linkage Logic}
Fraud syndicates operate by opening multiple accounts at the same physical bank branch using similar registered profiles. MuleTrace implements real-time entity resolution to discover these relationships:

\begin{equation}
\text{Link}(A, B) \iff (Branch_A = Branch_B) \land (Occupation_A = Occupation_B)
\end{equation}

This rule connects accounts sharing both \textbf{Bank Branch Code (\texttt{F3887})} and \textbf{Registered Occupation (\texttt{F3891})}, highlighting suspicious clusters.

\subsection{Node Visual threat Coding}
In the interactive dashboard network map, nodes are color-coded according to threat probability:
\begin{itemize}[itemsep=2pt]
    \item \textbf{Purple Diamond ($\blacklozenge$):} Simulated accounts created in Sandbox Mode.
    \item \textbf{Red Circle ($\bullet$):} Flagged high-risk accounts ($P \ge Threshold$).
    \item \textbf{Orange Circle ($\bullet$):} Warning status accounts ($20\% \le P < Threshold$).
    \item \textbf{Green Circle ($\bullet$):} Safe low-risk accounts ($P < 20\%$).
\end{itemize}

\subsection{Graph Focus Mode}
To isolate target networks in dense graphs, clicking a node activates Focus Mode. This interface action triggers a visual transformation:
\begin{itemize}
    \item Fades all unrelated nodes and edges in the graph to \textbf{15\% opacity}.
    \item Maintains full opacity on the selected node and its direct neighbors.
    \item Updates the side panel with matching profile features and local attributions.
\end{itemize}

This allows investigators to quickly trace the reach of money mule networks across branches.

% ── PAGE 10: XAI & NLG Engine ─────────────────────────────────────────────────
\newpage
\section{Explainable AI \& Dynamic NLG}

\subsection{Local Feature Attribution}
MuleTrace calculates Local Feature Attribution (LFA) to identify which variables have the greatest impact on an account's classification. The attribution value for feature $i$ is calculated relative to clean cohort medians:

\begin{equation}
\text{Attribution}_i = (X_i - \mu_i) \times \text{Importance}_i
\end{equation}

where $X_i$ is the feature value, $\mu_i$ is the median feature value for clean accounts, and $\text{Importance}_i$ is the feature importance calculated by the model. This isolates the specific anomalies that drove the classification.

\subsection{Natural Language Generation}
The Local Feature Attribution values are passed to the NLG engine, which maps the anonymized feature IDs to their semantic labels and generates compliance narratives:

\begin{lstlisting}[caption={Sample Generated Narrative Summary}]
AI Compliance Audit Verdict: SUSPECT MULE
This account is flagged as a SUSPECT MULE with an evaluated risk probability of 65.0%. The primary drivers triggering this alert are: Fraud Link Flag (1.0), Net Balance Ind. (381286.0), Account Type (Savings). Crucially, the entity resolution engine detected a direct shared device or branch linkage (Fraud Link) to another confirmed suspect mule account.
\end{lstlisting}

The system also displays specific sentences for each active driver, such as:
\begin{itemize}[itemsep=2pt]
    \item \textit{``Customer is registered as a student. Students are frequently targeted by criminal syndicates for money mule recruitment.''}
    \item \textit{``Customer is young (19 years old). Under 23 year-olds are the primary demographic targeted for student mule accounts.''}
\end{itemize}

Investigators can use these outputs to generate a pre-filled Suspicious Activity Report (SAR), reducing reporting time from hours to seconds.

% ── PAGE 11: Appendix ─────────────────────────────────────────────────────────
\newpage
\section*{Appendix: Technical Components}
\addcontentsline{toc}{section}{Appendix: Technical Components}

\renewcommand{\arraystretch}{1.3}
\begin{tabularx}{\textwidth}{>{\bfseries}p{2.5cm} >{\ttfamily\small}p{3.5cm} p{1.5cm} X}
\toprule
\textbf{Component} & \textbf{Source Files} & \textbf{Tool} & \textbf{Function} \\
\midrule
FastAPI Server & app.py & Python & Implements API endpoints and scoring logic. \\
ML Models & xgb\_hybrid\_model.json, lgb\_hybrid\_model.txt & XGBoost, LGBM & Bimodal classification engines. \\
Metadata Store & model\_metadata.json & JSON & Caches cohort medians and categorical mappings. \\
Frontend UI & static/index.html & HTML, JS & Renders dashboards and network visualizations. \\
Entity Resolution & static/index.html & vis.js & Handles network construction and Focus Mode. \\
NLG Engine & static/index.html & JS & Generates natural language compliance reports. \\
\bottomrule
\end{tabularx}

\vspace{16pt}
\noindent\textbf{Local Compilation & Running Commands:}

\begin{lstlisting}[language=bash, caption={Installation and local service execution}]
# Clone repository
git clone https://github.com/karthikeya3342/MuleTrace.git
cd MuleTrace

# Install dependencies
pip install -r requirements.txt

# Run FastAPI backend locally
uvicorn app:app --host 127.0.0.1 --port 8000 --reload
\end{lstlisting}

\vspace{12pt}
\noindent\textbf{Deployment Endpoints:}\\
- Deployed Live Application: \url{https://muletrace-ysh4.onrender.com/}\\
- GitHub Repository: \url{https://github.com/karthikeya3342/MuleTrace}

\end{document}
